\phantomsection\addcontentsline{toc}{chapter}{Abstract} % Don't remove me!
\chapter*{\centering{\Large Abstract}}

FPGA toolchains typically perform static compilation and mapping of fixed designs. In practice, audio signals and user-defined module graphs vary over time, presenting an opportunity to predictively reconfigure the FPGA fabric. This thesis proposes an operator-theoretic and analytic framework that models DSP modules as dynamical systems, using observables and new signal correlation metrics to forecast workload and adjust resource allocation. Unlike conventional flows, this approach can anticipate the evolving demands of a modular audio effect chain (including non-linear modules) and remap kernels (e.g. heavy multipliers or trigonometric units) onto LUTs, DSP slices, or BRAMs in accordance to resource demands. Expected contributions include: (1) an analytical model of audio DSP modules grounded in operator theory (Koopman embeddings) and state-space representations; (2) coherent metrics of signal interaction for resource tuning; and (3) a prototype architecture enabling runtime FPGA optimization in an accessible form. This work bridges a gap between control-theoretic modeling and FPGA architecture, offering audio engineers and embedded developers a more adaptive FPGA DSP platform that can improve throughput, support complex nonlinear \& stochastic audio effects and enable massive real-time parallelisation.
